\section{A successful extraction can still lose the fact}

Consider an Inline XBRL amount whose literal is \code{1,234}, whose admitted
numeric transform treats the comma as grouping, whose scale is $-3$, and whose
context and unit identify a particular dollar observation. The resulting amount
is $1234\times10^{-3}=1.234$. Scaling that normalized value again gives 0.001234
and is wrong. An adapter must retain the literal, transform, scale, context,
unit, nil status and original occurrence so that this error can be found.

The same principle applies to dates and dimensions. A lexical instant of
31 December remains distinct from the exclusive 1 January boundary used for
an interval. Publication and ingestion are separate timestamps. Two equal
numbers with different contexts remain separate occurrences. A nil fact is
unknown, not zero. Dropping a typed dimension or context reference yields a
preservation failure for that fact even if the resulting Markdown looks clear.

The inspected Arelle interfaces support native XBRL contexts, dimensions,
transformations and validation. Its internal object IDs do not replace stable
source identity. The inspected Docling source offers useful layout and table
proposals, but its examined XBRL graph path does not retain every required
fact/context/unit/typed-dimension field. That path needs a preservation sidecar
or must report unsupported conversion for affected records. These are source
audit findings, not measured accuracy rankings of the tools.

The retained Apple annual filing makes the dependency problem concrete. Its
issuer archive includes the instance, the issuer's schema and its calculation,
definition, label and presentation linkbases. Those files also refer to US GAAP
and SEC Reporting Taxonomy definitions, SEC schemas and XBRL definitions held
elsewhere. The issuer archive alone therefore cannot supply every definition
needed to interpret its tagged facts. For this particular archive, a static
inspection now resolves the recorded imports, link references and typed-domain
references to retained files, checks target namespaces and fragment identifiers,
and checks that instance facts have declared concepts, contexts and units.
This establishes a source dependency map. XML Schema and XBRL validation, the
issuer's interpretation of each concept, and the completeness of the wider
disclosure package remain separate questions.

An offline parser needs more than a map of intended filenames. At the inspected
Arelle source version, the cache lookup can return a local pathname when the
file is absent; subsequent loading records an \code{IOerror}. The proposed
adapter must check each mapped file and its hash before loading, retain the
missing dependency and affected records, and return an unresolved-dependency
result. It must not infer successful retrieval from a nonempty pathname or
silently substitute a newer taxonomy. Original schemas remain unchanged;
an external resolver or isolated cache associates their declared addresses
with the retained bytes. Any built-in cache fallback must belong to the same
checked source set \citep{arellepinned2026}.

Nor is a successful API return necessarily a validation verdict. In the same
source version, the controller's return value takes accumulated validation
errors into account only when \code{validationExitCode} is enabled. The proposed
adapter sets it explicitly, requests validation, disables automatic formula
execution, and examines the returned status, diagnostic messages and model
errors together. It keeps the model open while copying the necessary facts,
contexts, units and original occurrence references, then closes the session.
One session runs at a time in its owning thread within a process; the inspected
API uses shared package and plugin state. These are source-grounded interface
requirements for a later implementation, not results from an executed Arelle
adapter \citep{arellepinned2026}.

Each adapter therefore pins code, dependency, model and taxonomy versions,
licenses, input and output schemas, transformations and field-loss behavior.
Preparation preserves bytes and receipts before parsing, inventories the
package, attaches occurrences and contexts, records normalization, checks
required fields independently and submits candidates for review. Partial
conversion retains both successful units and the exact failures or unexplored
frontier. Conversion counts cannot accept a package or an accounting treatment.

The proposed extraction comparison uses the same retained files, roots, budgets
and taxonomy resources. Manual located source units and transparent native
extraction provide the reference; applicable pdfplumber/Arelle configurations,
plain Docling and explicit context/footnote repairs form the candidate ladder.
Transcription, context, table structure, mapping, source support, retention,
historical admission and cost are separate outcomes. Lost hashes or locators,
look-ahead, silent omissions and untracked downloads veto the affected exchange.
The current 31 development annotations cannot establish a heldout accuracy
ranking, and cell counts within one report are not independent observations.

Reuse also has distribution conditions. Arelle's Apache and Docling's MIT code
licenses do not by themselves clear every model, taxonomy or dataset. ACTUS
technical-specification reuse retains its share-alike obligations; QuantLib,
Z3 and ledger tools retain their own code and dependency terms. The Beancount
and hledger comparison needs an explicit distribution boundary for copyleft
components. License admission records actual versions, notices, models and
data, rather than treating a repository's headline license as universal.

The actual taxonomy notices show why that distinction changes the design.
The inspected FASB notices permit specified financial reporting and unchanged
incorporation in works that explain or assist implementation, subject to their
notice and third-party conditions. They do not provide a general permission to
edit taxonomy definitions. Keeping the original files and recording location
mappings separately respects that distinction. The notices retrieved for this
inspection are unversioned; their relationship to the terms attached to the
2024 files still needs to be established before a distribution decision.
XBRL's public-document notice separately directs software and schemas to a
Software Notice, so permission for an explanatory document cannot simply be
applied to a schema bundle. Exact dependency licences, component-specific terms
and the intended distribution remain explicit preparation tasks
\citep{fasbtaxonomyuses,fasbsrtuses,xbrlcopyright}.
